\documentclass[12pt,a4paper]{scrartcl}
\usepackage{rawfonts}
\usepackage{amssymb}

\usepackage{color}
\definecolor{astral}{RGB}{46,116,181}
\definecolor{myblue}{rgb}{0.25,0.25,0.5}

\usepackage{sectsty}
\sectionfont{\color{astral}}

\usepackage[colorlinks=true, urlcolor=blue, linkcolor=red]{hyperref}

\setlength{\parskip}{1.5ex plus 0.5ex minus 0.5ex}

\begin{document}
\title{\color{myblue}{Comparing K-Means and Other Algorithms for Data Clustering\break in Assembly. }}

\author{Héctor S. Enrique}
\date{ march 17, 2026}
\maketitle

\section{Purpose:}

	Some weeks ago I found the article: Comparing K-Means and Others Algorithms for Data Clustering by Nicolás Descartes. With C\# source code.

Look very interesting to translate that to Assembly. Still is a work in progress but look well.

There is some kind of abuse of Collections, because is very easy to write that in C\#. I removed most obvious exagerations (perhaps author try to make more clear the algorithm, not sure).

 \begin{itemize}
	 \item K-Means strategy need initial randomness, then probably you have to run several times before to find the better solution.
	 \item Hierarchical strategy need to define termination, and in these cases are number of clusters.
	 \item Density-based spatial clustering of applications with noise really benefit from Collections, and also need a couple of Sorted Vectors.
 \end{itemize}

Anyway is a good challenge, and still I have to collect some leaks

\section{Running the Program: }

    \begin{itemize}
        \item Step 1: Select from Menu Data one data set. If you want to use default set or already loaded set, then skip this step. 
        \item Step 2: Select from Menu File  $\displaystyle >$
 Load Data. If you want to use already loaded set, then skip this step. 
        \item Step 3:  Select from Menu Strategy the algorithm you want. If you want to use default stratey or already select algorithm, then skip this step. 
        \item Step 4 : Only if you seleted K-Means or Hierarchical strategies. Select from  Menu Cluster the number of clusters. In this examples the author use 3-clusters for data1 and data2, and  2-clusters for data 3. 
        \item Step 5 : Select from Menu File $\displaystyle >$ Run 
        \item Step 6 : Go to Step1 to test more things. 
     \end{itemize}

\section{Additional Reading}

	\href{https://www.codeproject.com/Articles/5375470/Comparing-K-Means-and-Others-Algorithms-for-Data-C}{Comparing K-Means and Others Algorithms for Data Clustering} by Nicolás Descartes 

\thispagestyle{empty}

\end{document}